The sky shown in the planetarium is stationary, and the observer is on the
surface of the Earth. Visible on the sky are: a comet, the Moon and a nova of
about 2nd magnitude. From the 11th minute, a grid representing horizontal
coordinates will be projected on the sky, and will remain on until the end of
the question. You are given a map of the sky for epoch J2000.0, using a polar
projection with a linear scale in declination, covering stars down to about
5th magnitude.

Note: Azimuth is counted from $0^\circ$ to $360^\circ$ starting at S through
W, N, E.

\begin{description}
\item[(A)] On the map of the sky, mark (with a cross) and label the nova
  (mark it ``N'') and the Moon (mark it with a Moon symbol) and draw the
  shape and position of the comet.
\item[(B)] In the table below, circle only those objects which are above the
  astronomical horizon. Note: you will lose 1 point for every incorrect
  answer.

  \begin{center}
  \begin{tabular}{lll}
  M20 -- Trifid Nebula & $o$ Cet -- Mira & $\delta$ CMa -- Wezen \\
  $\alpha$ Cyg -- Deneb & M57 -- Ring Nebula & $\beta$ Per -- Algol \\
  $\delta$ Cep -- Alrediph & $\alpha$ Boo -- Arcturus &
    M44 -- Praesepe (Beehive Cluster) \\
  \end{tabular}
  \end{center}
\item[(C)] When the coordinate grid is visible, mark on the map the northern
  part of the local meridian (from the zenith to the horizon) and the
  ecliptic north pole (with a cross and marked ``P'').
\item[(D)] For the displayed sky, give the:
  geographical latitude of the observer $\varphi$;
  Local Sidereal Time $\theta$;
  time of year, by circling the calendar month:
  Jan, Feb, Mar, Apr, May, Jun, Jul, Aug, Sep, Oct, Nov, Dec.
\item[(E)] Give the names of the objects whose approximate horizontal
  coordinates are:
  azimuth $A_1 = 45^\circ$ and altitude $h_1 = 58^\circ$;
  azimuth $A_2 = 278^\circ$ and altitude $h_2 = 20^\circ$.
  (If you can, use Bayer designations, IAU abbreviations and Messier numbers
  or English or Latin names.)
\item[(F)] Give the horizontal coordinates (azimuth, altitude) of:
  Sirius ($\alpha$ CMa): $A_3$, $h_3$;
  the Andromeda Galaxy (M31): $A_4$, $h_4$.
\item[(G)] Give the equatorial coordinates ($\alpha$, $\delta$) of the star
  marked on the sky with a red arrow.
\end{description}

% FIGURE NEEDED: polar-projection sky map (epoch J2000.0, stars to 5th
% magnitude) on which the marking tasks A, C and L are carried out
